\documentclass[11pt,a4paper]{article}

\usepackage[T1]{fontenc}
\usepackage[utf8]{inputenc}
\usepackage{lmodern}
\usepackage[margin=2.6cm]{geometry}
\usepackage{microtype}
\usepackage{amsmath,amssymb}
\usepackage{booktabs}
\usepackage{array}
\usepackage{enumitem}
\usepackage{titlesec}
\usepackage[hidelinks]{hyperref}
\usepackage{xcolor}

\definecolor{rulegrey}{gray}{0.45}
\titleformat{\section}{\large\bfseries}{\thesection}{0.7em}{}
\titleformat{\subsection}{\normalsize\bfseries}{\thesubsection}{0.7em}{}
\titlespacing*{\section}{0pt}{1.6em}{0.6em}
\titlespacing*{\subsection}{0pt}{1.2em}{0.4em}
\setlist[itemize]{leftmargin=1.2em,itemsep=0.25em,topsep=0.3em}
\setlist[enumerate]{leftmargin=1.5em,itemsep=0.45em,topsep=0.4em}
\newcommand{\hi}[1]{\textbf{#1}}
\newcolumntype{d}{>{\raggedleft\arraybackslash}p{1.45cm}}

\begin{document}

\begin{center}
{\LARGE\bfseries The EM chapters, re-examined}\\[0.45em]
{\large Confirmatory evidence for Chapters 8--11}\\[0.9em]
{\small Internal report for discussion with Jerome and Marc \quad$\cdot$\quad 16 September 2026}
\end{center}

\vspace{0.4em}
\hrule height 0.8pt
\vspace{1.2em}

\noindent
This report states what sixteen-seed confirmatory runs, executed on the Bocconi cluster over the
past two days, say about the chapters from EM onwards. Three of the four chapters survive intact;
one conclusion reverses; one comparison the thesis explicitly declared out of reach has now been
made. Section~\ref{sec:summary} is the short version, Section~\ref{sec:method} states what we are
doing and why it is trustworthy, Sections~\ref{sec:ch8}--\ref{sec:ch11} go chapter by chapter, and
Section~\ref{sec:ask} lists the five decisions that need taking.

\section{Summary}\label{sec:summary}

\begin{enumerate}
\item \hi{The capacity conclusion of Chapter~10 reverses.} It was an artefact of a stopping rule
that fired for the single-Gaussian arm and for no other. Under a protocol that selects the
checkpoint identically for every capacity, the single Gaussian sits exactly on the second-order
floor at every sample size ($0.167$) and every richer capacity goes far below it ($0.023$ at
$M=2048$), favouring the mixture in $16$ of $16$ seeds at every size.

\item \hi{The two rates of Chapter~9 are a property of the channel, not of the optimiser.} The EM
estimator attains the Cram\'er--Rao bound of the noised likelihood to within a factor $1.02$ on
average. The shape coordinate is not slow because EM is slow; the data stop carrying it. Its
information bound crosses the coordinate's own range at $t=0.98$, and that is precisely where the
fit collapses to a Gaussian innovation.

\item \hi{The comparison of Chapter~10 now runs at a matched information interface.} The thesis
attached a scope condition saying it could not do this. Given the same single noised view per
sequence that the estimator sees, the best network reaches $0.155$ against the estimator's
$0.023$: a paired factor $7.3$, interval $[6.2,8.7]$. The unmatched comparison, at $2.83\times$,
\emph{understates} the gap.

\item \hi{Locality is a ceiling, not only a prior.} The radius-4 window head is the best network at
$M=32$ and the worst of the five weight-shared architectures at $M=2048$, moving only
$0.121\rightarrow0.118$ as data quadruples. A bounded receptive field cannot represent a message
that has crossed the chain --- which is exactly what the chain's score is.

\item \hi{Score accuracy transfers to generation, and the second-order score fails there
visibly.} Driven through the same reverse SDE on identical Brownian paths, the fitted kernel
produces sequences whose innovation excess kurtosis is $1.887$ against the reference's $1.896$,
with a log-likelihood gap of $0.022$ under the true law. The covariance-matched score gives
$-0.008$ --- an essentially Gaussian chain --- at a gap of $1.71$, and the isotropic score a gap of
$46.2$.

\item \hi{A claim of ours was refuted by our own data, and is recorded as such.} ``The learned
networks sit at or above the second-order floor'' had been drafted into the new abstract. At
$M=2048$ five of six architectures sit below it. The defensible statement is about the information
interface, not about a barrier the networks cannot pass.
\end{enumerate}

\section{What we are doing, and why the numbers can be trusted}\label{sec:method}

Three changes to how the work is run account for why these results settle questions the thesis
left open.

\paragraph{One frozen protocol, one varied factor.}
Everything the arms share is now a single frozen object: the bundle sizes and their random
streams, the sizes $M\in\{32,128,512,2048\}$, the sixteen confirmatory seeds and four development
seeds, the EM iteration ladder and its initialisation, the network training recipe and screen
budget, the estimand, and the scoring regions. Each experiment varies exactly one declared factor
against that background, so a difference between two arms is attributable to that factor. Two arms
that differ in more than one thing are not compared; where a comparison genuinely needs two
factors to move, the intermediate arm is run as well.

\paragraph{Pairing by construction.}
Every row records the SHA-256 of the validation and test bundles it was scored on. Rows from
different experiments are joined on (seed, size, bundle hash), and a generator that is handed rows
with mismatched hashes refuses to emit a number. This is what makes ``EM--BP against the best
network'' a paired comparison across two experiments rather than two averages.

\paragraph{Frozen outputs.}
A result becomes citable only when its output directory is pinned: every CSV and provenance file
hashed into the registry entry together with the commit that produced it. After that, any byte
change fails the build, and the manuscript refuses to compile if a generated number disagrees with
the artefact it came from. A re-run is a new artefact under a new key, never an overwrite.
Confirmatory seeds are declared before the run and never used for selection; development seeds
never enter a reported number.

\paragraph{Amendments are recorded, not absorbed.}
Two design changes were made this week before the runs that used them, and both are in the
registry with timestamps: the EM ladder was extended from $400$ to $1200$ iterations because $400$
was binding at the largest sample size, and the largest strength of the site-varying
misspecification mechanism was removed because it is not a valid chain --- at $\rho=0.85$ its site
coefficients exceed $1$ and the implied covariance is not positive definite. No number from the
invalid point exists.

\section{Chapter 8 --- EM for the parameters: holds}\label{sec:ch8}

Nothing in the derivation needed revision. Fisher's identity still gives the exact gradient of the
discretised marginal log-likelihood from one forward--backward pass, and the finite-difference
check agrees to about $10^{-9}$. Two things were added.

The algebra is now machine-checked where checking is cheap. A Lean~4 module builds against Mathlib
with no \texttt{sorry} and no axiom, covering four constants: the Hermite attenuation for degrees
up to four, the fourth-cumulant difference under variance matching, the factor $\tfrac12$ in the
Gaussian information rate, and the coefficient $\kappa_4/6$ in the score gap. It verifies the
algebraic skeleton and not the analysis: the remainder bounds and the measure-level statements
remain informal, and the write-up says so in those words rather than claiming more.

The confirmatory ladder was extended to $1200$ iterations. At $M=2048$ the richer capacities were
still selecting the top of the $400$-iteration ladder; they still select the top of the extended
one in $10$ of $16$ seeds, so the errors reported for those arms in Table~\ref{tab:capacity}
remain upper bounds. That can only understate the effect described there.

\section{Chapter 9 --- the two rates: same fact, better reason}\label{sec:ch9}

The chapter explains the asymmetry through EM's fixed-point map: the Jacobian is the missing
information and the shape coordinate has the smaller eigenvalue, so the fitted innovation shape
settles at a median of $229$ updates against $80$ for the correlation. That measurement stands,
and the stopping-rule warning built on it stands with it.

What is new is that the asymmetry is no longer only a statement about the optimiser. Two results
bracket it.

\paragraph{The channel, from theory.}
Under the Ornstein--Uhlenbeck channel every Hermite mode of degree $m$ is attenuated exactly by
$e^{-mt}$, so the Fisher information one noised sequence carries about order-$m$ structure decays
as $e^{-2mt}$: correlation as $e^{-4t}$, non-Gaussian shape as $e^{-8t}$. Computed on the chain by
exact inference the measured slopes are $-3.97$ and $-8.01$ against the predicted $-4$ and $-8$,
and the constant in the score gap between the true score and that of the covariance-matched
Gaussian is reproduced to within one per cent.

\paragraph{The estimator, from measurement.}
Fitting the transition law from sequences observed at a single level $t$ ($M=512$, $16$ seeds, one
noised view per sequence), the EM estimator sits \emph{at} the Cram\'er--Rao bound of the noised
likelihood: a ratio of $1.02$ on average over the levels where either coordinate is informative,
never worse than $1.18$ for the correlation or $1.51$ for the fourth cumulant.

\medskip
\noindent
Taken together these say the chapter's phenomenon belongs to the channel. If the estimator is
efficient, the two rates cannot be blamed on EM. The bound for the shape coordinate grows as
$e^{4t}$ and crosses the coordinate's own range ($\kappa_4=3$) at $t=0.98$; the fit collapses to a
Gaussian innovation at exactly that level, while the correlation estimate is still within $0.009$
of the truth and survives past the crossover $t_*=1.17$.

\paragraph{One limitation, stated plainly.}
The corollary predicts errors growing as $e^{2t}$ and $e^{4t}$, and this run cannot test those
exponents. At $\rho=0.85$ the shape estimate saturates before the crossover, so the asymptotic
window does not exist at this configuration. Fitting slopes through a saturated estimate would
have returned $3.64$ and $0.09$ against predictions of $2$ and $4$; the analysis therefore reports
no slope, and a companion run at $\rho=0.5$ with $M=2048$ --- where the crossover falls at $0.55$
and the window is reachable --- is registered and queued.

\section{Chapter 10 --- the comparison: one reversal, one replacement, two new facts}\label{sec:ch10}

\subsection{The capacity sweep reverses}

The chapter concluded that no mixture capacity improves on a single Gaussian innovation, and
flagged its own caveat in the accompanying figure: the share of fits stopped at the iteration cap,
and with a narrowest component below the grid's resolution floor, both climbed steeply with $C$.
That caveat was the entire effect. In the archival sweep the likelihood tolerance fired in $31$ of
$32$ single-Gaussian cells at a median of $71$ iterations --- close to where a validation protocol
would stop --- and in none of the $96$ cells with $C\ge4$, which therefore reported checkpoints an
order of magnitude past their optimum. ``The single Gaussian is as good'' was ``the single Gaussian
is the only arm that stopped in the right place''.

Table~\ref{tab:capacity} re-runs the sweep with the checkpoint selected on a per-seed validation
bundle identically for every capacity.

\begin{table}[htbp]
\centering
\small
\caption{Capacity under equal stopping. Schedule-pooled relative score error on all sites, mean
over $16$ seeds, every arm on identical held-out bundles. The last column is the same capacity
fitted to the innovation law itself --- the $C$-component mixture closest in Kullback--Leibler
divergence --- an infinite-data reference for the family rather than a bound. The isotropic score
is at $0.669$ at every size.}
\label{tab:capacity}
\begin{tabular}{l dddd c d}
\toprule
& \multicolumn{4}{c}{sequences $M$} & & \\
\cmidrule(lr){2-5}
kernel & 32 & 128 & 512 & 2048 & & $M=\infty$ \\
\midrule
LMMSE floor      & 0.167 & 0.167 & 0.167 & 0.167 & & 0.167 \\
$C=1$            & 0.174 & 0.168 & 0.168 & 0.167 & & 0.167 \\
$C=2$            & 0.084 & 0.050 & 0.035 & 0.028 & & 0.039 \\
$C=4$            & 0.093 & 0.053 & 0.037 & 0.025 & & 0.0057 \\
$C=8$            & 0.099 & 0.055 & 0.036 & 0.023 & & 0.0022 \\
$C=16$           & 0.100 & 0.056 & 0.036 & 0.022 & & 0.0007 \\
\bottomrule
\end{tabular}
\end{table}

Three readings follow. The single-Gaussian kernel lands on the second-order floor at every sample
size, which is Proposition~1 read as a statement about an estimator rather than about a
projection. Every richer capacity goes below that floor, by a factor $7.9$ at the largest size; the
paired difference between $C=1$ and $C=8$ is $0.144$ with bootstrap interval $[0.141,0.147]$ and
favours the mixture in $16$ of $16$ seeds. And the best capacity rises with the data, from $C=2$ at
$M=32$ to $C=16$ at $M=2048$ --- the ordinary variance-to-bias trade, visible now that the stopping
rule no longer masks it.

\paragraph{The check that closes the loop.}
Two things changed at once between the thesis sweep and this one: the protocol, and the estimand
(held-out log-evidence per edge there, relative score error here). Two simultaneous changes are not
an explanation. A registered run now scores the \emph{same frozen fits} in the thesis's own
estimand, leaving the protocol as the only difference. A single-seed pass already shows the
reversal --- log-evidence per edge $-1.3116$ for $C=8$ against $-1.3146$ for $C=1$, with the true
kernel at $-1.3076$ --- and the sixteen-seed result is running.

\subsection{The data-efficiency comparison is replaced}

The chapter reports $6.6$--$10.9\times$ against unstructured denoising score matching and
$2.3$--$6.1\times$ against a baseline given locality and weight sharing, with three scope
conditions, the sharpest being that the arms are asymmetric in supplied information and that the
comparison therefore does not rank learning algorithms at a matched interface.

Six architectures were screened under one recipe --- $68$ configurations, four development seeds,
three restarts, $32{,}000$ steps each --- and the winner of each family was run on the sixteen
confirmatory seeds (Table~\ref{tab:baselines}).

\begin{table}[htbp]
\centering
\small
\caption{Under one training recipe, $16$ confirmatory seeds, identical held-out bundles. All-site
schedule-pooled relative score error. The upper block is trained on clean sequences with a fresh
corruption at every gradient step; the lower block sees what the estimator sees.}
\label{tab:baselines}
\begin{tabular}{l dd l}
\toprule
estimator & $M=32$ & $M=2048$ & information supplied \\
\midrule
fully connected              & 0.517 & 0.208 & clean, fresh corruption \\
window head ($r=4$)          & 0.186 & 0.118 & clean, fresh corruption \\
bidirectional GRU            & 0.201 & 0.060 & clean, fresh corruption \\
transformer                  & 0.231 & 0.063 & clean, fresh corruption \\
1-D U-Net                    & 0.227 & 0.066 & clean, fresh corruption \\
dilated TCN                  & 0.223 & 0.068 & clean, fresh corruption \\
\midrule
best matched-interface net   & ---   & 0.155 & one noised view \\
EM--BP ($C=8$)               & 0.099 & 0.023 & one noised view \\
\midrule
\textit{LMMSE floor}         & \textit{0.167} & \textit{0.167} & --- \\
\bottomrule
\end{tabular}
\end{table}

Two numbers replace the thesis's ranges. Against the best of six architectures under the standard
recipe, the estimator is ahead by a paired factor $2.83$, interval $[2.37,3.42]$, in $16$ of $16$
seeds at $M=2048$, and by $2.07$ at $M=32$. Against the best network trained at the estimator's own
interface --- one noised view per sequence, no clean chain, checkpoint selected on the implicit
objective of a held-out noisy bundle and never on the true score --- the factor is $7.3$, interval
$[6.2,8.7]$. Implicit score matching on that bundle, which uses no clean target at all, lands at
$0.165$.

The second number is the one the thesis said was out of reach, and it is the stronger claim: the
gap is not an artefact of giving the networks less, because when they are given the same, the gap
widens.

It also follows that one sentence drafted for the new paper is false and has been withdrawn.
``Every network sits at or above the second-order floor'' holds only at the smallest sample size;
at $M=2048$ five of the six sit below it. Given enough clean sequences and fresh corruptions a
neural denoiser plainly does learn fourth-order structure. The honest statement concerns the
information interface.

\subsection{Two facts the thesis does not contain}

\paragraph{Locality is a ceiling.}
The window head is the best network at $M=32$ and the worst of the five weight-shared
architectures at $M=2048$, moving only from $0.121$ to $0.118$ as the data quadruples, against a
standard error of $0.001$, while the global architectures continue to $0.060$. A radius-4 window
cannot represent a message that has travelled the length of the chain, and the chain's exact score
is precisely such a message. This connects the architecture comparison to the bulk-posterior
analysis the chapter already contains.

\paragraph{Length extrapolation separates the arms cleanly.}
Fitted at $L=32$ and evaluated against the exact score of longer chains
(Table~\ref{tab:length}), the fitted kernel is flat: the law is local, so inference at any length
is the same exact recursion. The convolutional and recurrent networks improve slightly, as they
should, since longer chains have proportionally fewer boundary sites. The transformer with learned
relative positions degrades by an order of magnitude --- a failure of the positional encoding
beyond its training length, not of the score it learned.

\begin{table}[htbp]
\centering
\small
\caption{Length extrapolation at $M=2048$: pooled relative score error of estimators fitted at
$L=32$ and evaluated against the tree-exact score at longer chains, mean over $16$ seeds.}
\label{tab:length}
\begin{tabular}{l dddd}
\toprule
estimator & $L=32$ & $L=64$ & $L=128$ & $L=256$ \\
\midrule
EM--BP ($C=8$)     & 0.023 & 0.023 & 0.023 & 0.023 \\
bidirectional GRU  & 0.060 & 0.057 & 0.057 & 0.056 \\
1-D U-Net          & 0.066 & 0.064 & 0.062 & 0.062 \\
dilated TCN        & 0.068 & 0.066 & 0.065 & 0.064 \\
window head        & 0.118 & 0.099 & 0.088 & 0.082 \\
transformer        & 0.063 & 0.243 & 0.514 & 0.703 \\
\bottomrule
\end{tabular}
\end{table}

\section{Chapter 11 --- which answers change}\label{sec:ch11}

\paragraph{RQ5, recovering an unknown transition.}
The answer holds and gains a sentence. The E-step is structurally exact, recovery is demonstrated
numerically over the studied regime, and the uneven settling of the coordinates is real. We can now
add that the estimator is statistically efficient, so the unevenness belongs to the channel and not
to the algorithm. The stopping-rule warning becomes sharper rather than weaker: a rule read off the
correlation trace certifies a kernel whose shape coordinate the data may no longer identify at all.

\paragraph{RQ6, what supplying the Markov assumption is worth.}
The answer changes in its numbers and in its scope. The ranges $6.6$--$10.9\times$ and
$2.3$--$6.1\times$ are superseded by $2.83\times$ against the best of six architectures under one
recipe. The third scope condition --- that the comparison does not rank methods at a matched
information interface --- can be removed, because that comparison now exists and gives
$7.3\times$. The capacity subsection of the same chapter must be rewritten in the opposite
direction.

\paragraph{The retired claim.}
The capacity result should be stated as retired rather than quietly dropped, together with the
diagnosis of why it arose. The same applies to the claim of ours that our own screen refuted this
week. The project keeps a negative-results record for exactly this, and both entries are in it.

\section{Generation and the boundary of the advantage}\label{sec:gen}

Two results completed while this report was being written. Both were pre-registered and both
address questions Chapter~10 raises but does not settle.

\paragraph{Does score accuracy transfer to generation?}
Every arm drives the same reverse SDE from the same initialisation along the same Brownian path, so
the comparison is paired within a seed and the integrator is common to all of them. The
discretisation was measured separately and is not what separates the arms: against an $800$-step
reference the strong error is $0.0190$, $0.0087$ and $0.0033$ at $100$, $200$ and $400$ steps,
roughly halving per doubling, while the excess kurtosis of the generated sequences is stable at
$2.00$--$2.05$ across the whole ladder.

\begin{table}[htbp]
\centering
\small
\caption{Generation at $M=2048$, sixteen seeds, every arm on the same Brownian path within a seed.
The deficit is the paired difference in mean $\log p_{t_{\min}}$ per sequence against the arm driven
by the exact score, in nats, with a seed bootstrap; the kurtosis column is $|$generated $-$
reference$|$, the reference sample having excess kurtosis $1.936$.}
\label{tab:generation}
\begin{tabular}{l l dd}
\toprule
score driving the sampler & deficit vs exact (nats) & \multicolumn{1}{c}{kurt.\ error} & \multicolumn{1}{c}{cov.\ error} \\
\midrule
exact (tree BP)         & ---                        & 0.119 & 0.070 \\
EM--BP ($C=8$), fitted  & $+0.09$\ \ $[+0.00,+0.19]$ & 0.177 & 0.077 \\
bidirectional GRU       & $+0.32$\ \ $[+0.17,+0.46]$ & 0.385 & 0.102 \\
window head             & $+1.05$\ \ $[+0.79,+1.31]$ & 0.474 & 0.290 \\
LMMSE (covariance twin) & $+1.70$\ \ $[+1.68,+1.71]$ & 1.939 & 0.069 \\
isotropic               & $+46.2$                    & 1.936 & 0.908 \\
\bottomrule
\end{tabular}
\end{table}

The ordering is the one the score errors predict, and the fitted kernel is nearest the exact score
by a wide margin: its deficit is $0.09$ nats against $0.32$ for the best network and $1.70$ for the
covariance twin. The deficit is small but real --- an earlier four-seed run suggested the fitted
score matched the exact one, and sixteen seeds show it does not. The second-order score is the
instructive failure. On the covariance error it is the best arm in the
table, marginally ahead of the exact score, which is what Proposition~1 predicts. On the fourth
cumulant it is the worst: it generates an essentially Gaussian chain with the right covariance. A
generative model evaluated only on second-order distances would score it a success.

\paragraph{Where does the advantage stop?}
The structured estimator assumes a first-order Markov chain. Four Gaussian departures from that
assumption were swept, each with a scalar strength $\varepsilon$ and an exact score of its own, so
every arm is scored against the truth of the perturbed model rather than against the chain.

\begin{table}[htbp]
\centering
\small
\caption{Misspecification: ratio of the window head's pooled error to EM--BP's, at zero strength
and at the largest strength swept. Above $1$ the misspecified chain estimator still wins; the last
column interpolates where the ratio crosses $1$. Eight seeds, sizes $512$ and $2048$ pooled.}
\label{tab:misspec}
\begin{tabular}{l l d d d l}
\toprule
mechanism & what it breaks & \multicolumn{1}{c}{$\varepsilon=0$} & \multicolumn{1}{c}{$\varepsilon_{\max}$} & \multicolumn{1}{c}{ratio there} & break-even \\
\midrule
rank-one global latent   & nothing structural & 10.5 & 1.00 & 3.96 & not reached \\
AR(2) second coefficient & the order          & 10.9 & 0.40 & 0.97 & $\approx0.40$ \\
long-range coupling      & the precision band & 7.9  & 0.40 & 1.00 & $\approx0.40$ \\
site-varying $\rho$      & homogeneity        & 9.8  & 0.14 & 0.99 & $\approx0.14$ \\
\bottomrule
\end{tabular}
\end{table}

The pattern is clean and interpretable. A rank-one global latent is absorbed by the fitted
correlation and never reverses the ordering, because it does not break the conditional
independence structure the estimator exploits. The three mechanisms that do break it --- adding an
order, widening the precision band, or making the chain inhomogeneous --- each reverse it at a
measurable strength, and the site-varying mechanism reverses it soonest. This gives the chapter the
boundary it currently lacks: the advantage is not a general property of ``having structure'', it is
specific to the structure being approximately right, and how wrong it can be is now quantified.

\section{Status of the evidence, and what is still running}\label{sec:status}

\begin{table}[htbp]
\centering
\small
\caption{Confirmatory runs. \emph{Frozen} means every file in the output directory is pinned by
SHA-256 in its registry entry, with the commit that produced it; a generator refuses to read an
output that is not frozen or whose bytes have changed.}
\label{tab:status}
\begin{tabular}{l l l}
\toprule
run & what it establishes & state \\
\midrule
E01  & order-by-order decay verified on the chain              & frozen \\
E02a & capacity under equal stopping                            & frozen \\
E02b & capacity against innovation shape at fixed covariance    & frozen \\
E03  & six architectures, matched interface, length             & frozen \\
E06  & learning through noise, against the information bound    & frozen \\
E07  & the second-order floor across correlation                & frozen \\
E09  & infinite-data reference per capacity                     & frozen \\
\midrule
E10  & capacity in the thesis's own estimand                    & frozen \\
E04  & grid refinement on the fitted models                     & running \\
E05  & generation under each score, paired noise                & frozen \\
E08  & misspecification and its break-even                      & frozen \\
E06b & the asymptotic exponents at $\rho=0.5$                   & frozen, inconclusive \\
\bottomrule
\end{tabular}
\end{table}

\emph{Refinement} is now complete and is as clean as this kind of check gets. Over the whole ladder
--- $401$, $801$ and $1601$ grid points, half-widths $8$ and $10$ --- every arm's error agrees to
five decimal places and the reference field itself moves by at most $10^{-5}$; the worst seed's
spread in the headline number is zero at that precision. No number in this report is an artefact of
the quadrature at the frozen configuration. Whether the same holds at the other correlations,
families and chain lengths the project uses is a separate question, and a registered sweep is now
running to answer it rather than assume it.

\paragraph{Three infrastructure faults, for the record.}
Jobs were inheriting the login shell's $600$-second CPU limit, which killed one run outright and
left another holding an idle compute node for ten hours while the scheduler reported it running.
Two drivers wrote their results only at the very end, so a timeout would have lost everything; both
now write per cell and resume. And the site-varying misspecification mechanism is undefined at its
largest registered strength. All three are fixed; the third is a registered amendment.

\section{Decisions requested}\label{sec:ask}

\begin{enumerate}
\item \hi{How should the capacity reversal be presented?} Either rewrite the capacity section of
Chapter~10 around the new result and retain the old figure as a documented artefact, or keep both
and make the diagnosis itself the finding. The second is more honest about how the error surfaced
and costs roughly a page.

\item \hi{Is the matched-interface comparison the headline, or a robustness check?} It is the
stronger claim and it removes a scope condition the thesis carried. Promoting it changes the
framing from ``structure beats networks'' to ``structure beats networks at equal information, and
the unmatched comparison understates it''.

\item \hi{Where does the efficiency result belong?} The estimator attaining the Cram\'er--Rao bound
is what turns the two rates from an observation about EM into a statement about information. It
reads naturally in Chapter~9 or in the theory chapter, and should be stated once with a reference
from the other.

\item \hi{How far should the Lean verification go?} Four algebraic constants are machine-checked.
The remainder bounds are analysis and would be a substantial effort; the projection operator behind
Proposition~1 is the one load-bearing step still in prose. Is the current scope the right stopping
point?

\item \hi{Do we pursue the asymptotic exponents, or claim the rates only where they are
established?} The run at $\rho=0.5$ has reported and the window it opens is three levels wide, too
narrow to pin an exponent (Section~\ref{sec:ch9}). The window grows only logarithmically in $M$, so
closing it properly means a much smaller $\rho$ and a much larger $M$, or switching to an estimand
that does not saturate --- which is a different claim. The alternative, which is what the write-up
currently does, is to establish the rates exactly by computation and the estimator's efficiency by
measurement, and to state the sample-complexity consequence only where the coordinate is still
identifiable.
\item \hi{Does the misspecification boundary belong in Chapter~10 or in the conclusions?} It is
the cleanest statement we have about when the method should \emph{not} be used --- a rank-one
global latent never reverses the ordering, while breaking the order, the precision band or the
homogeneity each reverse it at a measurable strength. It reads as a limitation, but it is a
measured one, which is unusual enough to be worth foregrounding.
\end{enumerate}

\vspace{1.4em}
\hrule height 0.4pt
\vspace{0.7em}
\noindent
{\footnotesize Every number above comes from a registered experiment whose output directory is
pinned by SHA-256 in its registry entry, together with the commit that produced it and the seeds it
used. The protocol shared by all arms is a single frozen object, so two arms differ only in the one
factor each experiment varies. Runs were executed on the Bocconi HPC. Supporting material: the
claim--evidence matrix, the negative-results record, the experiment registry and the
reproducibility notes in the project repository.}

\end{document}
